% =============================================================================
% Chapter 6 -- Iteration 1: Design and Implementation  (owners: all three)
% Grounded in the repository as it stands at commit 671f3a2 (9 Oct 2026).
% Everything described as implemented below was verified by running the
% build/typecheck/lint commands quoted in Section "Demonstration Evidence".
% Backend, ingestion, model-server and contracts work does NOT exist yet and is
% stated as such -- see Section "Iteration Outcome" and Chapter 8 defects.
% =============================================================================
\chapter{Iteration 1: Design and Implementation}
\label{ch:iteration1}

This chapter describes what has actually been built. File names, route paths,
component names and command outputs in this chapter come from the repository
itself; where a capability is simulated rather than implemented, it is labelled
\emph{mocked} at the point of description rather than in a footnote.

% ---------------------------------------------------------------------------
\section{Iteration Objectives}
\label{sec:iter1-objectives}

The iteration objective recorded in Chapter~3 was to lay the foundations every
later component depends on: a reproducible repository scaffold, the shared
contracts written as code, one event ingested end to end, deterministic links
resolved on a fixture, and a measurement of local inference on the hardware
actually available.

Work completed in this iteration was redirected by the mid-evaluation
requirements. The objectives actually pursued were:

\begin{enumerate}
  \item Freeze the requirements and design baseline so that implementation has
        something to be checked against (Chapters~4 and~5).
  \item Build the reviewer surface screen-for-screen against the approved
        Stitch design, with the domain vocabulary encoded in types rather than
        in prose, so that the interface can be demonstrated and defended before
        the back end exists.
  \item Produce the twelve design diagrams as committed, editable artefacts.
  \item Construct and validate the first package of evaluation data, because the
        panel required the dataset to be validated rather than named.
\end{enumerate}

Objective 1 and objectives 3--4 were met. Objective 2 was met for the interface
only. The original backend objectives \emph{were not started}; they are carried
forward unchanged into the plan in Chapter~10.

% ---------------------------------------------------------------------------
\section{Selected Use Cases / Features}
\label{sec:iter1-features}

The iteration implements the presentation path of three use cases from
Chapter~4 and defers their data paths:

\begin{longtable}{p{0.10\textwidth} p{0.34\textwidth} p{0.24\textwidth} p{0.26\textwidth}}
\caption{Use cases selected for iteration 1, with the acceptance criteria each one must satisfy.}\label{tab:iter1-usecases}\\
\hline\textbf{Use case} & \textbf{What the iteration delivers} & \textbf{Status} &
  \textbf{What is still missing} \\
\hline\endfirsthead
\hline\textbf{Use case} & \textbf{What the iteration delivers} & \textbf{Status} &
  \textbf{What is still missing} \\
\hline\endhead
UC-01 & The change view: selecting a pull request shows its linked criteria,
  each with a verdict, citations, reason and missing evidence; the verification
  passes are inspectable.
  & Implemented against typed mock data
  & Live ingestion, real packets, a real verifier run \\
UC-02 & Link review: inferred and unconfirmed links are listed with their
  provenance and can be confirmed or rejected in the interface.
  & Implemented against typed mock data
  & A graph that stores the decision, an API that persists it \\
UC-03 & Finding lifecycle: acknowledge, resolve, suppress and reopen with a
  reason, rendered with an audit trail.
  & Implemented against typed mock data
  & Server-side enforcement and the audit table \\
UC-04 to UC-06 & Requirement divergence, pre-filled configuration, repository
  eligibility scoring.
  & Not started (policy and analyser do not exist yet)
  & Entirely back end \\
\hline\end{longtable}

% ---------------------------------------------------------------------------
\section{Design}
\label{sec:iter1-design}

\subsection{Design source and tokens}

The interface was not designed twice. Every screen was exported from the Stitch
project \emph{Unified Dev Collaboration Hub} into
\texttt{docs/stitch-reference/screens/}, which holds 20 HTML exports with 19
PNG renders. The application's theme in
\texttt{apps/web/src/app/globals.css} is a one-to-one port of the token block
found inside those exports: the \texttt{@theme} block defines the surface,
primary, secondary, tertiary, error and outline colours, the type scale, the
spacing scale and the radii. Screens are checked against their exported source,
and ad-hoc hex values are forbidden in components.

Fonts are self-hosted through \texttt{next/font} (Inter and JetBrains Mono) with
the Material Symbols icon set served from \texttt{/fonts/}. Nothing is fetched
from a content distribution network, which is what allows the app to run fully
offline and keeps the local-first commitment of Chapter~1 true even for the
prototype.

\subsection{Layout and information architecture}

\texttt{src/components/layout/app-shell.tsx} provides the frame used by every
route: a 240-pixel side rail and a sticky top bar. Navigation is declared
centrally in \texttt{src/lib/nav.ts} as two arrays, \texttt{primaryNav} (seven
entries: dashboard, conflict inbox, conflict alerts, graph explorer, activity
feed, link management, audit trail) and \texttt{workspaceNav} (nine entries:
policy configuration, policy simulation, search, workspace admin,
notifications, Slack delivery, onboarding, evaluation, study runner), with
\texttt{isActive()} deciding the highlighted item. One declaration means a new
screen appears in the rail once, in the right group.

\subsection{Domain vocabulary as types}

The most consequential design decision in this iteration is that the frozen
label taxonomy of Chapter~4 exists as code. \texttt{src/lib/types.ts} declares
\texttt{VERDICTS} as a four-element const tuple
(\texttt{SUPPORTED}, \texttt{CONTRADICTED}, \texttt{PARTIAL},
\texttt{INSUFFICIENT\_EVIDENCE}) from which the \texttt{Verdict} type is
derived, plus \texttt{FreshnessState} (\texttt{CURRENT}/\texttt{STALE}/
\texttt{UNKNOWN}), a closed \texttt{ReasonCode} union
(\texttt{SOURCE\_STALE}, \texttt{CITATION\_UNRESOLVED},
\texttt{LINK\_NOT\_ELIGIBLE}, \texttt{BELOW\_THRESHOLD},
\texttt{SCHEMA\_VIOLATION}, \texttt{ABSTAINED}), and separate
\texttt{CriterionResult} fields for verdict, freshness, review flag and reason.
Freshness and review-required are therefore structurally prevented from being
folded into the verdict, which is the rule the dataset protocol also relies on.
The same file defines \texttt{Citation} with a \texttt{locator}, an
\texttt{excerpt}, a source \texttt{tool} and a \texttt{resolved} flag.

% ---------------------------------------------------------------------------
\section{Implementation}
\label{sec:iter1-implementation}

\subsection{Reviewer surface (Abdullah)}

\texttt{apps/web} is a Next.js App Router application. Nineteen routes exist
(one folder per screen under \texttt{src/app/}), and the front end totals 90
TypeScript and TSX files of 23{,}564 lines: 57 component files contributing
17{,}010 lines, and an eleven-module mock layer of 3{,}163 lines.

The screens with the most behaviour are:

\begin{itemize}
  \item \textbf{Conflict detail} (\texttt{src/app/conflicts/[id]/page.tsx}) --
        the evidence case. It renders the criterion results, explanation box,
        affected artifacts, evidence chain, verification passes, audit trail and
        resolution dock from
        \texttt{src/components/conflict-detail/}. It exports
        \texttt{generateStaticParams()}, which pre-renders one page per mock
        finding, so eight detail pages are generated statically at build time.
        Severity and status are mapped through the typed
        \texttt{severityMeta} and \texttt{statusLabel} records rather than
        free-text strings.
  \item \textbf{Conflict inbox} (\texttt{/conflicts}) -- the findings table with
        stat cards, a multi-tier toolbar and severity/status filtering.
  \item \textbf{Graph explorer} (\texttt{/graph} and
        \texttt{/graph/artifacts}) -- node and edge views whose coordinates
        share a single 0--100 space so the two views agree.
  \item \textbf{Link management} (\texttt{/links}) -- the suggestion stack,
        link wizard, provenance modal and directory, where an
        \texttt{inferred} or \texttt{unconfirmed} suggestion can be confirmed or
        rejected.
  \item \textbf{Policy workbench and simulation} (\texttt{/policies},
        \texttt{/policies/simulation}) -- rule definitions, the workflow state
        map, and a simulation lab that reports why a rule did or did not fire.
\end{itemize}

Shared UI lives in \texttt{src/components/ui/} (\texttt{Icon}, \texttt{Panel},
\texttt{Badge}, \texttt{ToolPill}, \texttt{Button}, \texttt{SegmentedTabs},
\texttt{StatCard}, \texttt{PageHeader}), so every screen uses the same primitives
for the same concepts.

\subsection{Requirements and design artefacts (Mishal)}

Chapter~4 was written against the proposal sections listed in its header, and
introduced FR-22 to FR-29 in agreement with the team (traceability anchors,
symbol cards, comment--code consistency, repository eligibility, the two
confirmation thresholds and pre-filled configuration). Four TikZ diagrams
(\texttt{usecase.tex}, \texttt{domain-model.tex}, \texttt{erd.tex},
\texttt{anchor-reanchoring.tex}) plus the shared style file
\texttt{unify-diagram-styles.tex} and \texttt{build.sh} were committed in
\texttt{717a7af} and wired into Chapters~4 and~5. The ERD carries thirteen
tables whose fields match the schema sections of Chapter~5.

\subsection{Evaluation data package (Mishal construction, Haaniah annotation)}

The pilot dataset package in \texttt{docs/dataset/} was constructed in this
iteration: 8 repositories screened with accept/reject/defer reasons
(\texttt{01\_repository\_catalogue.csv}), a frozen case schema
(\texttt{02\_case\_schema.json}), a 66-line annotation manual with worked
examples, 15 case packets with head patches and test logs, two annotation
sheets, an adjudication log template, a proposed repository-level split
manifest, an integrity report and a dataset card draft. Three scripts live under
\texttt{docs/dataset/tools/}: \texttt{build\_cases.py},
\texttt{check\_integrity.py} and \texttt{agreement.py}.

Haaniah's contribution in this iteration is the annotation itself: sheet
\texttt{annotator\_B\_haaniah.csv} is filled for all 15 cases with verdict,
freshness, review flag, reason code, evidence spans, justification and time per
case, committed in \texttt{8657955}. Her three diagrams (ingestion sequence,
state machines, deployment) were committed in \texttt{ae0a03c}, and the GenAI
register was maintained in \texttt{a5cdf50}.

\subsection{What does not exist yet}

Stated plainly, because the worksheet requires the boundary to be visible:

\begin{itemize}
  \item There is no \texttt{apps/api}, \texttt{apps/worker},
        \texttt{packages/contracts}, \texttt{db/migrations} or
        \texttt{tools/mock-vendor}; there is no \texttt{docker-compose.yml} and
        no CI configuration in the repository.
  \item No event has ever been ingested: there is no adapter code, no
        normaliser, no checkpoint and no graph in the current tree.
  \item The context-packet builder, the verifier loop, the verdict gate and the
        local model server do not exist. No model has been downloaded, and no
        latency, memory or quantisation measurement has been taken.
  \item The front end reads from \texttt{src/lib/mock/}. Nothing in
        \texttt{src/} performs a network request.
\end{itemize}

The pre-defence proof-of-concept engine that used to live in
\texttt{unify\_poc/} (an ingest $\rightarrow$ normalise $\rightarrow$ link
$\rightarrow$ rule $\rightarrow$ notify pipeline with a Flask UI) was deleted
from the tree by the \texttt{066ef56} ``start fresh'' commit. Two of its files
remain, and they no longer form a runnable artefact: see the defect recorded in
Chapter~8.

% ---------------------------------------------------------------------------
\section{Completed Features / Module}
\label{sec:iter1-completed}

\begin{longtable}{p{0.30\textwidth} p{0.16\textwidth} p{0.14\textwidth} p{0.30\textwidth}}
\caption{Modules and features completed during iteration 1, and their state in the repository.}\label{tab:iter1-completed}\\
\hline\textbf{Capability} & \textbf{State} & \textbf{Evidence} & \textbf{Substitute
  and how it will be removed} \\
\hline\endfirsthead
\hline\textbf{Capability} & \textbf{State} & \textbf{Evidence} & \textbf{Substitute
  and how it will be removed} \\
\hline\endhead
Nineteen-route reviewer surface, all routes rendering
  & Real & \texttt{npm run build} prerenders 29 pages; \texttt{7b411b1}
  & None needed \\
Design tokens ported from the exported Stitch screens
  & Real & \texttt{globals.css} \texttt{@theme} block; \texttt{b09e05a}
  & None needed \\
Typed domain vocabulary (verdicts, freshness, reason codes, citations)
  & Real & \texttt{src/lib/types.ts}; typecheck passes
  & None needed \\
Finding lifecycle, link confirmation, policy simulation screens
  & Real UI, \textbf{mocked data}
  & \texttt{src/lib/mock/} (11 modules, 3{,}163 lines)
  & Replace \texttt{src/lib/mock} imports with an API client; the mock modules
    are shaped like the future response objects \\
Requirements, FR/NFR, acceptance and traceability tables
  & Real & Chapter~4; \texttt{33b00e7} & None needed \\
Twelve diagrams, one build script
  & Real & \texttt{diagrams/} + \texttt{rendered/}; \texttt{717a7af},
    \texttt{ae0a03c}, \texttt{d2d7d0c}
  & Commit stamp per diagram still outstanding \\
Evaluation-data pilot (15 cases, manual, scripts, integrity report)
  & Real & \texttt{docs/dataset/}; \texttt{66e85c7}
  & None needed \\
Second-annotator agreement statistics
  & In progress & Annotator B complete, annotator A pending
  & Not a substitute -- a pending measurement (Ch.~8) \\
Event ingestion, normalisation, idempotency, checkpoints
  & \textbf{Not started} & -- & -- \\
Evidence graph, linking, rules, incremental dispatch
  & \textbf{Not started} & -- & -- \\
Context packet, verifier, verdict gate, model server
  & \textbf{Not started} & -- & -- \\
\hline\end{longtable}

% ---------------------------------------------------------------------------
\section{Technical Implementation Details}
\label{sec:iter1-details}

\paragraph{Static generation of the detail route.}
\texttt{conflicts/[id]} is the only dynamic route. Its
\texttt{generateStaticParams()} maps over the mock \texttt{findings} array, so
the build emits one HTML file per finding. This is how the conflict detail page
can be opened during a demonstration without a server, and it is also the reason
the build output lists eight \texttt{CONF-*} paths.

\paragraph{Typed status rendering.}
In the same route, \texttt{severityMeta} maps each \texttt{Severity} to a chip
label, tone and left-bar class, and \texttt{statusLabel} maps each
\texttt{FindingStatus} to a fixed string. Because both maps are typed
\texttt{Record<Severity, ...>}, adding a status without updating the map is a
compile error rather than a blank label at run time.

\paragraph{Mock layer shape.}
\texttt{src/lib/mock/findings.ts} exports \texttt{findings} and
\texttt{getFinding(id)}; \texttt{conflict-detail.ts} exports helpers such as
\texttt{formatUtcFull}, \texttt{getPolicy} and \texttt{packetHash}. The
\texttt{packetHash} helper exists because the design in Chapter~5 requires the
evidence case to name the packet it came from; in the mock it is a stable
synthetic hash, and it becomes a real digest when the packet builder lands.
\texttt{link-suggestions.ts} documents at the top of the file that every
suggestion it returns is \texttt{inferred} / \texttt{unconfirmed} and is
therefore listable and confirmable but never a decisive finding.

\paragraph{Diagram build.}
\texttt{diagrams/build.sh} renders all twelve diagrams in one run: fragment
sources are wrapped in a \texttt{standalone} preamble that inputs
\texttt{unify-diagram-styles}, and Haaniah's three self-contained documents are
compiled as they are, at a different resolution to reproduce her original
renders. Output lands in \texttt{diagrams/rendered/}.

% ---------------------------------------------------------------------------
\section{Algorithms and Technical Decisions}
\label{sec:iter1-algorithms}

\begin{enumerate}
  \item \textbf{Type-derived contracts before contracts exist.} With
        \texttt{packages/contracts} not yet created, the verdict vocabulary was
        frozen in TypeScript first so that screens, mock data and the future
        schema compile against one definition. When the Zod contracts land, the
        types move with them; nothing else in the front end needs to change
        because components consume \texttt{Verdict} and
        \texttt{CriterionResult}, not string literals.
  \item \textbf{Mock data as a shaped module, not inline constants.} Every
        screen imports from \texttt{src/lib/mock/} rather than declaring its own
        arrays, so the substitution is a single import path per module. This is a
        deliberate choice to keep the implemented/mock boundary removable.
  \item \textbf{Static rendering where the data is static.} Routes are static by
        default, which makes the demo independent of a back end and makes
        regressions visible as build failures.
  \item \textbf{Coordinate normalisation in the graph view.} Node and edge
        positions are stored in one shared 0--100 space instead of two, so the
        explorer and the artifact view cannot drift apart visually
        (\texttt{d535476}).
  \item \textbf{No runtime image or font fetching.} Avatars are initials,
        thumbnails are gradients, fonts are local files. This is a privacy and
        offline decision as much as a performance one.
\end{enumerate}

% ---------------------------------------------------------------------------
\section{Tools, Frameworks and Technologies}
\label{sec:iter1-tools}

\begin{longtable}{p{0.26\textwidth} p{0.16\textwidth} p{0.50\textwidth}}
\caption{Tools, frameworks and technologies used in iteration 1, and the reason each was chosen.}\label{tab:iter1-tools}\\
\hline\textbf{Tool} & \textbf{Version used} & \textbf{Role in this iteration} \\
\hline\endfirsthead
\hline\textbf{Tool} & \textbf{Version used} & \textbf{Role in this iteration} \\
\hline\endhead
Next.js & 16.4.0 & App Router front end, static generation \\
React & 19.3.0 (declared) & Component model \\
TypeScript & 5.9.3 (declared) & Domain types; \texttt{tsc --noEmit} gate \\
Tailwind CSS & 4.3.3 (declared) & \texttt{@theme} token block, utility classes \\
ESLint & 9 (declared) & Lint gate via \texttt{eslint-config-next} \\
Node.js & v24.14.1 & Toolchain used for the verified build \\
Google Stitch & MCP-integrated project & Screen design and export \\
Python & 3.12 & Dataset construction, integrity and agreement scripts \\
pytest & 9.1.1 & Test runner for the surviving proof-of-concept file \\
TikZ / \texttt{tikz-uml} & -- & The twelve diagrams \\
\hline\end{longtable}

Declared versions are those in \texttt{apps/web/package.json}; the Node version
is the one the build reported.

% ---------------------------------------------------------------------------
\section{Demonstration Evidence}
\label{sec:iter1-demo}

The following commands were run from \texttt{apps/web} on the reporting commit
and their results are reproduced verbatim:

\begin{verbatim}
$ npm run typecheck        # tsc --noEmit
  (exit code 0, no diagnostics)

$ npm run lint             # eslint .
  (exit code 0, no findings)

$ npm run build            # next build
  Running TypeScript ... Finished TypeScript in 3.2s ...
  Generating static pages ... (29/29) in 1075ms
  Route (app): 19 routes listed, /conflicts/[id] emits
  CONF-1046, CONF-1047, CONF-1048 and five more paths
\end{verbatim}

\noindent
What a demonstration can show today: the full reviewer journey through mock
data --- select the conflict inbox, open a conflict, read each criterion's
verdict and citations, inspect the verification passes, confirm or reject a
suggested link, and record a resolution --- entirely offline, from a production
build. What it cannot show today is any of that data arriving from a real tool,
because the ingestion, graph and verification layers described in Chapter~5 do
not exist yet. The boundary above is the honest one to state before a live demo.

% ---------------------------------------------------------------------------
\section{Iteration Outcome}
\label{sec:iter1-outcome}

The iteration produced a stable, buildable reviewer surface with the project's
domain vocabulary compiled into types, a complete requirements and design
baseline, twelve committed diagrams, and the first validated package of
evaluation data. It did not produce the ingestion, linking, verification or
model-serving foundations that the iteration plan in Chapter~3 named as its
deliverables, and it did not produce the technical spike that was meant to
retire the hardware risk early.

The consequences are stated rather than softened: the live demo at the mid
evaluation can only exercise the front end against mock data; the reliability
and verification claims in Chapters~4 and~5 remain unimplemented requirements;
and the risk that local inference does not fit the available hardware is still
open, with no measurement taken. The recovery plan, the resequenced work and the
scope consequences are in Chapter~10, and the current defect list is in
Chapter~8.
